\section{Forelesning 7: \emph{Implementering av IT-strategier}}
\label{sec:lec07}
\settermlabel{terms:lec07}

\subsection*{Oversikt}
Forelesning 7 (Ulrik Røhl) flytter fokus fra strategiformulering til
\emph{utførelse}. En IT-strategi som ikke implementeres, eksisterer ikke.
Forelesningen dekker Millers (1997) suksessfaktorer, Tawse \& Tabesh sitt
implementeringsrammeverk, Kotters 8-trinns endringsmodell, Scheins
kulturmodell, samt typiske årsaker til at IT-prosjekter feiler.
Bjørn-Andersen \& Hedman (2016) om Coops historiske IT-kompleksitet er det
sentrale empiriske bakteppet.

\subsection{Hovedteori}

\subsubsection*{Strategi vs.\ implementering}
\keypoint{IT-strategi må implementeres -- ikke bare formuleres.} Selv den beste
strategien feiler hvis utførelsen ikke følger. Implementering er ikke en
mekanisk fase etter formulering; den er en strategisk disiplin i seg selv.

\subsubsection*{Miller (1997): Fire (og fem) suksessfaktorer}
Miller definerer hva \guillemotleft{}vellykket implementering\guillemotright{} betyr gjennom fire faktorer:
\begin{itemize}
  \item \term{Completion} -- ble initiativet faktisk ferdigstilt?
  \item \term{Achievement} -- ble de planlagte målene oppnådd?
  \item \term{Acceptability} -- aksepterer interne og eksterne interessenter
    resultatet?
  \item \term{Propitiousness} -- gunstige omstendigheter rundt
    implementeringen.
\end{itemize}
Den utvidede sjekklisten legger til fire kontekstuelle faktorer:
\begin{itemize}
  \item \term{Backing} -- sterk og synlig støtte fra ledelsen.
  \item \term{Assessability} -- mulig å måle fremdrift og resultater
    underveis.
  \item \term{Specificity} -- konkrete, presise mål og tiltak.
  \item \term{Cultural receptivity} -- en organisasjonskultur som er
    mottakelig for endringen.
\end{itemize}

\begin{figure}[H]
\centering
\includegraphics[width=0.8\linewidth]{lec07-miller-four-success-factors.png}
\caption{Miller (1997) -- de fire grunnleggende suksessfaktorene for
implementering.}
\end{figure}

\begin{figure}[H]
\centering
\includegraphics[width=0.8\linewidth]{lec07-miller-five-implementation-factors.png}
\caption{Millers fem implementeringsfaktorer -- Backing, Assessability,
Specificity, Cultural receptivity, Propitiousness -- som en sjekkliste.}
\end{figure}

\subsubsection*{Tawse \& Tabesh (2021): Integrert rammeverk}
Et tre-laget rammeverk:
\begin{enumerate}
  \item \term{Conditions} -- \emph{Competency}, \emph{Commitment},
    \emph{Coordination} -- bestemmer Strategy Implementation Effectiveness.
  \item \term{Managerial Actions} -- \emph{Structural} og \emph{Interpersonal}
    handlinger som produserer betingelsene over.
  \item \term{Dynamic Managerial Capabilities} -- \emph{Cognition}, \emph{Social
    Capital}, \emph{Human Capital} -- som gjør riktig kombinasjon mulig.
\end{enumerate}
\keypoint{Tawse \& Tabesh er forelesningens primære analytiske linse for \guillemotleft{}hvordan
skal strategien faktisk implementeres?\guillemotright{}}

\begin{figure}[H]
\centering
\includegraphics[width=0.95\linewidth]{lec07-tawse-tabesh-implementation-framework.png}
\caption{Tawse \& Tabesh (2021) -- integrert implementeringsrammeverk i tre lag.}
\end{figure}

\subsubsection*{Kotter (2015): 8 trinn for å lede endring}
En klassisk, foreskrivende top-down-modell:
\begin{enumerate}
  \item Forøg presset (\emph{sense of urgency}).
  \item Etablér ledende koalisjon.
  \item Få visjonen og strategien klar.
  \item Kommuniser visjonen.
  \item Gjør det! (myndiggjør handlekraft).
  \item Skap synlige resultater (kortsiktige seire).
  \item Konsolider forbedringer.
  \item Implementer nye prosedyrer i kulturen.
\end{enumerate}
Brukes som kontrast til Tawse \& Tabesh sin mer emergente, kapabilitets-
baserte tilnærming.

\begin{figure}[H]
\centering
\includegraphics[width=0.9\linewidth]{lec07-kotter-eight-step-leading-change.png}
\caption{Kotters 8-trinnsmodell for å lede endring.}
\end{figure}

\subsubsection*{Schein \& Brown (1998): Kultur som lag}
Schein og Browns / Johnson et al.\ sin nested-circles-modell viser kultur som
fire konsentriske lag: \emph{Paradigm} (taken-for-granted antagelser) $\rightarrow$
\emph{Behaviours} $\rightarrow$ \emph{Beliefs} $\rightarrow$ \emph{Values}.
\keypoint{En strategisk plan kan være riktig, men feile fordi den krasjer mot
paradigm-laget -- antagelser som er så grunnleggende at de er usynlige.}

\begin{figure}[H]
\centering
\includegraphics[width=0.75\linewidth]{lec07-schein-culture-model.png}
\caption{Schein \& Brown -- kulturmodellens fire lag.}
\end{figure}

\subsubsection*{Hvorfor IT-prosjekter feiler}
Vanlige årsaker:
\begin{itemize}
  \item Gap mellom strategiformulering og utførelse -- de som planlegger
    snakker ikke med de som bygger.
  \item Mangel på kapabiliteter eller ressurser (teknisk, organisatorisk).
  \item Undervurdering av endringsbehovet -- IT er sosioteknisk, ikke bare
    teknisk.
  \item Teknisk kompleksitet og \term{legacy debt}.
  \item Feiltilpassede insentiver -- prosjektledere belønnes for å levere
    \guillemotleft{}på tid\guillemotright{} selv om verdien ikke realiseres.
  \item Svak \term{governance} -- uklare beslutningsrettigheter.
  \item Kulturelt motstand -- paradigm-laget i Schein-modellen.
\end{itemize}

\subsubsection*{IT-legacy-gjeld}

\begin{definition}{IT legacy debt}
Akkumulert kompleksitet fra historiske teknologivalg som begrenser fremtidig
fleksibilitet. Analogi til finansiell gjeld: rente betales i form av redusert
endringshastighet, økte vedlikeholdskostnader og koblingsavhengigheter.
\end{definition}

Bjørn-Andersen \& Hedman (2016) dokumenterte at Coop hadde betydelig
historisk IT-kompleksitet -- mange systemer, mange integrasjoner, lange
historiske valg som begrenset strategiske handlingsrom. Implikasjonen:
\keypoint{enhver ny digital strategi må aktivt styre kompleksitet; å bare
legge til nye funksjoner øker gjelden.}

\subsubsection*{Governance av IT-strategi}
\begin{itemize}
  \item \term{IT governance} -- rammeverket av beslutningsrettigheter og
    ansvarlighet som styrer IT-bruk.
  \item Sentralisert vs.\ desentralisert styring -- avveiing mellom kontroll
    og lokal fleksibilitet.
  \item \term{Hybrid governance} -- bevisst mix: noen beslutninger sentralt
    (plattform, sikkerhet), noen lokalt (anvendelse, brukeropplevelse).
  \item \term{Stage-gate} / pilotering -- testing før storskala-utrulling.
\end{itemize}

\subsection{Sentrale fagbegreper}\label{terms:lec07}
\begin{itemize}
  \item \term{Strategy Implementation Effectiveness}.
  \item \term{Conditions, Managerial Actions, Dynamic Managerial Capabilities}
    (Tawse \& Tabesh).
  \item \term{Sense of urgency}, \term{guiding coalition} (Kotter).
  \item \term{IT governance} -- beslutningsrettigheter og ansvarlighet.
  \item \term{Legacy debt} -- akkumulert teknisk kompleksitet.
  \item \term{Hybrid governance} -- mix av sentralisert og desentralisert.
  \item \term{Change management} -- den menneskelige siden av endring.
  \item \term{Stage-gate} -- beslutningspunkter mellom faser.
  \item \term{Benefits realization} -- disiplinen å sikre at lovet verdi
    faktisk leveres.
  \item \term{Paradigm} -- Scheins dypeste kulturlag.
\end{itemize}

\subsection{Implementeringsdisiplin: en sjekkliste}
\begin{enumerate}
  \item Er strategiens forutsetninger eksplisitte og testbare?
  \item Hvem eier hvilke beslutninger? (governance)
  \item Hvilke kapabiliteter trengs -- og finnes de? (Tawse \& Tabesh)
  \item Hvordan måles realisert verdi (ikke bare aktivitet)?
  \item Hvilken legacy-kompleksitet vil initiativet legge til eller redusere?
  \item Hvordan testes valgene før storskala-forpliktelse? (pilot/stage-gate)
  \item Hva i kulturen vil motstå dette? (Schein-paradigmet)
\end{enumerate}

\subsection{Eksempler og caser}
\begin{eksempelboks}{Coop (Bjørn-Andersen \& Hedman 2016)}
Historisk IT-kompleksitet og fragmentering på tvers av kjeder. Begrenset
strategiske valg og gjorde plattformkonsolidering kostbar. Implementering må
adressere både teknisk gjeld og organisatorisk kapabilitet -- ikke kun den
ene.
\end{eksempelboks}

\begin{eksempelboks}{Klassiske IT-fiaskoer}
Store ERP-implementeringer som overskrider budsjett og leverer mindre enn
lovet -- typiske eksempler på svikt i Millers \emph{specificity},
\emph{backing} eller \emph{cultural receptivity}, eller på manglende
\emph{commitment} hos Tawse \& Tabesh.
\end{eksempelboks}

\begin{eksempelboks}{Hybrid governance for Coop App}
Coop-first mandat, målbare KPI-er knyttet til butikkøkonomi, pilotering før
utrulling -- et konkret eksempel på hvordan implementeringsdisiplin
oversettes til en strategisk anbefaling.
\end{eksempelboks}

\subsection{Eksamensrelevans}
\begin{itemize}
  \item Bruke Tawse \& Tabesh som hovedlinse for implementeringsspørsmål.
  \item Forklare vanlige årsaker til at IT-strategier feiler i utførelse,
    med Miller (1997) som sjekkliste.
  \item Definere legacy debt og diskutere strategiske implikasjoner.
  \item Forklare hybrid governance og hvorfor det passer en omnichannel-
    aktør med blandet portefølje (jf.\ \autoref{sec:lec03}).
  \item Knytte implementering til pilotering: en pilot er en test av
    strategiens antagelser (jf.\ \autoref{sec:lec01}).
  \item Bruke Bjørn-Andersen \& Hedman (2016) som konkret evidens om Coop.
  \item Bruke Schein-modellen til å forklare kulturell motstand.
  \item Kontrastere Kotter (foreskrivende, top-down) mot Tawse \& Tabesh
    (kapabilitets-basert, emergent).
\end{itemize}
